% GENERATED FILE. Edit canonical lessons, then run npm run generate:books.
% canonical-source-hash: fnv1a64:fa73e04331a0efb9
% canonical-lessons: AR-C506-jafn, AR-C506-rimsh, AR-C506-litha, AR-C506-dhaqn, AR-C506-asab

\chapter{The face up close}
\label{ch:ar-the-face-up-close}
\hlchaptermodality{506}

\begin{chapteropening}
\textbf{By the end of this chapter:} \emph{I can name an eyelid, an eyelash, the gums, the chin and a nerve.}

Complete the last lesson of chapter 506: I can name an eyelid, an eyelash, the gums, the chin and a nerve.
\end{chapteropening}

\section[\texorpdfstring{\ar{جفن} (jafn)}{جفن (jafn)}]{\ar{جفن} (jafn) — an eyelid}

\label{lesson:AR-C506-jafn}



\begin{quote}
Before the new one: say the Arabic for a hip, then the Arabic for a joint.
\end{quote}

\subsection*{You'll want to know: \ar{جفن}}
\textbf{\ar{جفن}} — \emph{jafn} — \textquotedblleft{}an eyelid\textquotedblright{}.

\begin{quote}
\ar{فتح} \ar{المريض} \ar{جفنيه} \ar{ببطء}.
\end{quote}

\begin{quote}
\emph{fataḥa al-marīḍ jafnayhi bi-buṭʾ.}
\end{quote}

\begin{quote}
The patient slowly opened his eyelids.
\end{quote}

Plural: \ar{جفون}, \emph{jufūn}.

\begin{cousinweb}
\textbf{\ar{جفن}} (\emph{jafn}) is built on the native Arabic root \emph{j-f-n}, \textquotedblleft{}a sheath, a covering\textquotedblright{} (Lane). The eyelid is the eye's cover.
\end{cousinweb}

\begin{grammarlens}[title={What you've built}]
One more word for describing what you lived and felt.
\end{grammarlens}

\subsection*{Your turn}
\begin{itemize}
\raggedright
  \item \emph{Say it:} \emph{jafn}
  \item \emph{Say it:} \emph{jafn}, once more
  \item \emph{Say it:} \emph{fataḥa al-marīḍ jafnayhi bi-buṭʾ.} — The patient slowly opened his eyelids.
  \item \emph{Recall:} say the Arabic for a hip, then the Arabic for a joint, then say \emph{jafn} again
\end{itemize}

\begin{tcolorbox}[breakable,colback=teal!4,colframe=teal!35!black,title={Before you move on}]
What does \ar{جفن} mean? (\textquotedblleft{}An eyelid\textquotedblright{}.) Say it once more.
\end{tcolorbox}
\section[\texorpdfstring{\ar{رمش} (rimsh)}{رمش (rimsh)}]{\ar{رمش} (rimsh) — an eyelash}

\label{lesson:AR-C506-rimsh}



\begin{quote}
Before the new one: say the Arabic for a joint, then the Arabic for an eyelid.
\end{quote}

\subsection*{You'll want to know: \ar{رمش}}
\textbf{\ar{رمش}} — \emph{rimsh} — \textquotedblleft{}an eyelash\textquotedblright{}.

\begin{quote}
\ar{للطفلة} \ar{رموش} \ar{طويلة}.
\end{quote}

\begin{quote}
\emph{li-ṭ-ṭifla rumūsh ṭawīla.}
\end{quote}

\begin{quote}
The little girl has long eyelashes.
\end{quote}

Plural: \ar{رموش}, \emph{rumūsh}.

\begin{cousinweb}
\textbf{\ar{رمش}} (\emph{rimsh}) is a native Arabic word for an eyelash, usually traced to the root \emph{r-m-sh}; the origin of the sense is not settled. Its plural \emph{rumūsh} is the common form.
\end{cousinweb}

\begin{grammarlens}[title={What you've built}]
One more word for describing what you lived and felt.
\end{grammarlens}

\subsection*{Your turn}
\begin{itemize}
\raggedright
  \item \emph{Say it:} \emph{rimsh}
  \item \emph{Say it:} \emph{rimsh}, once more
  \item \emph{Say it:} \emph{li-ṭ-ṭifla rumūsh ṭawīla.} — The little girl has long eyelashes.
  \item \emph{Recall:} say the Arabic for a joint, then the Arabic for an eyelid, then say \emph{rimsh} again
\end{itemize}

\begin{tcolorbox}[breakable,colback=teal!4,colframe=teal!35!black,title={Before you move on}]
What does \ar{رمش} mean? (\textquotedblleft{}An eyelash\textquotedblright{}.) Say it once more.
\end{tcolorbox}
\section[\texorpdfstring{\ar{لثة} (litha)}{لثة (litha)}]{\ar{لثة} (litha) — the gums}

\label{lesson:AR-C506-litha}



\begin{quote}
Before the new one: say the Arabic for an eyelid, then the Arabic for an eyelash.
\end{quote}

\subsection*{You'll want to know: \ar{لثة}}
\textbf{\ar{لثة}} — \emph{litha} — \textquotedblleft{}the gums\textquotedblright{}.

\begin{quote}
\ar{قال} \ar{طبيب} \ar{الأسنان} \ar{إن} \ar{لثتي} \ar{سليمة}.
\end{quote}

\begin{quote}
\emph{qāla ṭabīb al-asnān inna lithatī salīma.}
\end{quote}

\begin{quote}
The dentist said that my gums are healthy.
\end{quote}

It is one word for all the gums.

\begin{cousinweb}
\textbf{\ar{لثة}} (\emph{litha}) is a native Arabic word for the gums. Lane notes that some file it under \emph{l-w-th}, from the flesh wound round the roots of the teeth; its root is not settled.
\end{cousinweb}

\begin{grammarlens}[title={What you've built}]
One more word for describing what you lived and felt.
\end{grammarlens}

\subsection*{Your turn}
\begin{itemize}
\raggedright
  \item \emph{Say it:} \emph{litha}
  \item \emph{Say it:} \emph{litha}, once more
  \item \emph{Say it:} \emph{qāla ṭabīb al-asnān inna lithatī salīma.} — The dentist said that my gums are healthy.
  \item \emph{Recall:} say the Arabic for an eyelid, then the Arabic for an eyelash, then say \emph{litha} again
\end{itemize}

\begin{tcolorbox}[breakable,colback=teal!4,colframe=teal!35!black,title={Before you move on}]
What does \ar{لثة} mean? (\textquotedblleft{}The gums\textquotedblright{}.) Say it once more.
\end{tcolorbox}
\section[\texorpdfstring{\ar{ذقن} (dhaqan)}{ذقن (dhaqan)}]{\ar{ذقن} (dhaqan) — a chin}

\label{lesson:AR-C506-dhaqn}



\begin{quote}
Before the new one: say the Arabic for an eyelash, then the Arabic for the gums.
\end{quote}

\subsection*{You'll want to know: \ar{ذقن}}
\textbf{\ar{ذقن}} — \emph{dhaqan} — \textquotedblleft{}a chin\textquotedblright{}.

\begin{quote}
\ar{وضع} \ar{الرجل} \ar{يده} \ar{تحت} \ar{ذقنه}.
\end{quote}

\begin{quote}
\emph{waḍaʿa ar-rajul yadahu taḥta dhaqanihi.}
\end{quote}

\begin{quote}
The man put his hand under his chin.
\end{quote}

Plural: \ar{ذقون}, \emph{dhuqūn}.

\begin{cousinweb}
\textbf{\ar{ذقن}} (\emph{dhaqan}) is built on the native Arabic root \emph{dh-q-n}, \textquotedblleft{}the chin\textquotedblright{} (Lane). Lane gives two plurals, \emph{adhqān} and \emph{dhuqūn}.
\end{cousinweb}

\begin{grammarlens}[title={What you've built}]
One more word for describing what you lived and felt.
\end{grammarlens}

\subsection*{Your turn}
\begin{itemize}
\raggedright
  \item \emph{Say it:} \emph{dhaqan}
  \item \emph{Say it:} \emph{dhaqan}, once more
  \item \emph{Say it:} \emph{waḍaʿa ar-rajul yadahu taḥta dhaqanihi.} — The man put his hand under his chin.
  \item \emph{Recall:} say the Arabic for an eyelash, then the Arabic for the gums, then say \emph{dhaqan} again
\end{itemize}

\begin{tcolorbox}[breakable,colback=teal!4,colframe=teal!35!black,title={Before you move on}]
What does \ar{ذقن} mean? (\textquotedblleft{}A chin\textquotedblright{}.) Say it once more.
\end{tcolorbox}
\section[\texorpdfstring{\ar{عصب} (ʿaṣab)}{عصب (ʿaṣab)}]{\ar{عصب} (ʿaṣab) — a nerve}

\label{lesson:AR-C506-asab}



\begin{quote}
Before the new one: say the Arabic for the gums, then the Arabic for a chin.
\end{quote}

\subsection*{You'll want to know: \ar{عصب}}
\textbf{\ar{عصب}} — \emph{ʿaṣab} — \textquotedblleft{}a nerve\textquotedblright{}.

\begin{quote}
\ar{أصاب} \ar{الحادث} \ar{عصبا} \ar{في} \ar{يده}.
\end{quote}

\begin{quote}
\emph{aṣāba al-ḥādith ʿaṣaban fī yadihi.}
\end{quote}

\begin{quote}
The accident damaged a nerve in his hand.
\end{quote}

Plural: \ar{أعصاب}, \emph{aʿṣāb}.

\ar{أعصاب} also means nerves as in calm: \ar{أعصابه} \ar{قوية}.

\begin{cousinweb}
\textbf{\ar{عصب}} (\emph{ʿaṣab}) is built on the native Arabic root \emph{ʿ-ṣ-b}, \textquotedblleft{}to bind, to tie tight\textquotedblright{} (Lane). In the old language ʿaṣab was a sinew, the cord that binds the body.
\end{cousinweb}

\begin{grammarlens}[title={What you've built}]
One more word for describing what you lived and felt.
\end{grammarlens}

\subsection*{Your turn}
\begin{itemize}
\raggedright
  \item \emph{Say it:} \emph{ʿaṣab}
  \item \emph{Say it:} \emph{ʿaṣab}, once more
  \item \emph{Say it:} \emph{aṣāba al-ḥādith ʿaṣaban fī yadihi.} — The accident damaged a nerve in his hand.
  \item \emph{Recall:} say the Arabic for the gums, then the Arabic for a chin, then say \emph{ʿaṣab} again
\end{itemize}

\begin{tcolorbox}[breakable,colback=teal!4,colframe=teal!35!black,title={Before you move on}]
What does \ar{عصب} mean? (\textquotedblleft{}A nerve\textquotedblright{}.) Say it once more.
\end{tcolorbox}
